%--------------------
% Packages
% -------------------
\documentclass{article}
\usepackage{fontspec}
\setmainfont{Times New Roman}
\setlength{\parindent}{0pt}
%-----------------------
% Begin document
%-----------------------
\begin{document} %All text i dokumentet hamnar mellan dessa taggar, allt ovanför är formatering av dokumentet
\title{Advanced programming methodologies for photonic processors}
\author{zhenyun xie}
\maketitle

\section{Introduction}

\subsection*{Motivation}
Talk about limitations of electrical integration circuit and challenges in the post-moore's law era:
\begin{itemize}
    \item Explosive growth in data transmission demand due to AI, ML..., Electronic processor reaches it's limitations, slow moore low, power, bandwidth, I/O bottle neck
    \item Requirements of new paradigm AI data center: energy efficiency, low latency, interconnect bandwidth, dynamic...
\end{itemize}

Talk about the state of the art to introduce OCS.
These leading industry demonstrate that OCS become a strategic technology choice for overcoming performance and energy efficiency bottlenecks in AI data centers.
\begin{itemize}
    \item google palomar, spine replacement, dynamic network
    \item nvidia paper, failure resilience
\end{itemize}

One of the core technic for OCS is use silicon photonic, talk about photonic processor:
\begin{itemize}
    \item Photonic integrated circuits are solutions, high bandwidth, low latency, passive, low power consumption
    \item Optical interconnects are solutions for AI data transmission, especially optical circuit switch, talk about advantages, network reconfiguration
\end{itemize}

Talk about the importance of photonic hardware programming:
\begin{itemize}
    \item Photonic integrated circuits they are programmable (PIP through MZI)
    \item Design and control complexity of the large-scale photonic processor
    \item Need a high level control layer to manage and configure photonic processor based products in datacenter
    \item Missing an application oriented programming abstractions for photonic processor design and controlling
\end{itemize}

\subsection*{Objectives}
Research objectives:
\begin{itemize}
    \item Provide an advanced, efficient, and universal programming methodology for programmable photonic processors in AI data centers
    \item Resolve how to efficiently map the datacenter applications to the photonic processor
    \item Address how to comprehensively design and simulate large scale OCS for datacenter
    \item Explore and develop an SDN of the layer 1 OCS for AI datacenter applications controller
\end{itemize}

Key contributions:
\begin{itemize}
    \item Graph-based advanced programming methodologies for photonic applications
    \item Framework for Optical Circuit Switch Architecture Design
    \item SDN control system for Optical Circuit Switch
\end{itemize}

\subsection*{Structure of the thesis}
Chapter 2: Theoretical and Programming Methodologies of Photonic Integrated Circuits
\begin{itemize}
    \item talk about phase shifter / MMI / MZI
    \item circuits, feedforward / circular
    \item graphic theory
    \item sdn theory
\end{itemize}

Chapter 3: Optical Networking Applications Based on Programmable Integrated Photonic
\begin{itemize}
    \item interconnect (result)
    \item beamsplitter (algorithm + result)
    \item switch (algorithm + result)
\end{itemize}

Chapter 4: optical circuit switch architectures for datacenter applications.
\begin{itemize}
    \item type of topologies Benes/Dilated Benes/Piloss/Dilated banyan/Double Dilated Benes
    \item blocking study packets traffic / ocs traffic
    \item sax simulations (ILs / SNR)
    \item VPI link simulation with a STD SOA specification?
\end{itemize}

Chapter 5: software defined network for OCS
\begin{itemize}
    \item custom yang model
    \item netconf
    \item gnmi
\end{itemize}

Chapter 6: discussions
\begin{itemize}
    \item address large-scale ocs (losses / crossing) problem
    \item ocs datacenter use cases maturity (reconfigurable topologies challenging to implement)
    \item missing an AI network optimization control software
    \item sdn standardization process
    \item conclusions
    \item future work
          \begin{itemize}
              \item OpenConfig
              \item new architectures
              \item multicore approaches
          \end{itemize}
\end{itemize}
\end{document}